\begin{abstract}
End-to-end driving policies are ranked by proxy evaluators---open-loop displacement, non-reactive
and two-stage pseudo-simulation---whose rankings have so far only been compared with each other,
never with a known answer.
We construct that answer: graded degradations of pretrained policies whose driving must get worse
with the degradation level, admitted as ground truth only after a closed-loop reference confirms the
order on one set of scenes, and then used to measure every evaluator on disjoint scenes.
Because sensitivity to harmful change can be achieved by penalising \emph{any} change, we add the
converse test: benign perturbations that alter only the part of a plan that a replanning vehicle never
executes, which a valid evaluator must ignore.
Across four policy families with up to three seeds, five evaluators and over three hundred evaluation runs, we
find that
(i)~only \num{RetainedCount} of six popular degradation families (\num{RetainedFamilies}) yield a
ground truth, even in closed loop;
(ii)~NAVSIM v1 and v2 penalise changes to the never-executed part of the plan by up to
\num{SpecPctNavsimVone} and \num{SpecPctNavsimVtwo} of the penalty for the matching harmful change,
while closed loop moves by less than half a point;
(iii)~the evaluator most sensitive to known failures, open-loop displacement, is also non-specific, so
sensitivity alone does not certify an evaluator; and
(iv)~on held-out scenes NAVSIM and closed loop rank real policies in opposite directions
(Kendall $\tau={}$\num{TauNavsimBridgesim}), including a state-of-the-art diffusion planner that NAVSIM
places \num{ContrastDDLTFNavsimVoneAbs} points above a simpler baseline and closed loop
\num{ContrastDDLTFBridgesimAbs} points below it.
We release the harness, all per-scenario scores and the statistical protocol.
\end{abstract}
